
This section assesses the feasibility of the ARGUS mission. 
To do so, we need to calculate certain key parameters essential to successful mission execution. 
These calculations should also validate our various parameter choices, 
particularly those relating to the selected LASER sources, 
the debris material and the target orbit.  
\subsection{Theoretical Calculation of Selected LASER Parameters}
\begin{itemize}
    \item \textbf{Ablation Threshold and Optimal Fluence (F)}\\
    The fluence of our LASER at a wavelength of $\lambda = 1064 nm$ with pulses of \textbf{nanosecond duration (ns)} 
    represents the energy deposited by this laser per unit area. It is expressed in \textbf{$J/cm^{2}$}. 
    The scientific literature places the theoretical ablation threshold between $F_{th} \approx 1 J/cm^{2}$ and $F_{th} \approx 2   J/cm^{2}$ (Philippe class reference).
    However, to obtain the required optimal propulsive effect in our case, we need to reach the maximum coupling fluence $F_{couplage}$ defined by: \\
    \begin{equation}
    \boxed{F_{couplage} \approx  2.5 F_{th} \approx  5 J/cm^{2}}
    \end{equation}
    for $F_{th} = 2 J/cm^{2}$\\

    \item \textbf{Pulse energy \textbf{$E_{P}$} and peak power \textbf{$P_{c}$}}\\
    Pulse energy represents the total energy contained in a single LASER pulse. 
    It is calculated using the following formula: 
    \begin{equation}
        \boxed{E_{P} = S *  F_{couplage} }
    \end{equation}
    $S$ represents the area illuminated by the LASER pulse. 
    For a selected focal spot diameter of $d = 0.3 cm$ and assuming a spherical spot, we obtain the value of $S$ :  
    \begin{equation}
      \boxed{  S = \frac{\pi d^{2}}{4} =\frac{\pi * (0.3)^{2}}{4} \approx 0.071 cm^{2} }
    \end{equation}
    With this value of $S$, we finally obtain: 
    \begin{equation}
        \boxed{E_{P} = 0.071 *5 = 0.35J }
    \end{equation}
    Thus, with this value of $E_{P}$, we can determine the instantaneous peak power \textbf{$P_{c}$} generated by our LASER. This power represents the energy flow per unit time. It is the maximum power reached by the LASER within the pulse itself. It is obtained by dividing the pulse energy by the pulse duration. 
    \begin{equation}
        \boxed{P_{c} = \frac{E_{P}}{\tau} = \frac{0.35}{5.10^{-9}} = 70MW}
    \end{equation}
    \\
    \item \textbf{Repetition rate $f_{rep}$ and average power $P_{moy}$}\\
    To accommodate the payload within a $6U$ CubeSat, we assume an average power of $P_{moy} = 1000 W$. 
    This value is not arbitrary; it is primarily dictated by the thermal dissipation limit of a $6U$ CubeSat \cite{Bekey1997,Phipps2012}.\\
    With this value, the pulse repetition rate is calculated using the formula below: \\
    \begin{equation}
        \boxed{f_{rep} = \frac{P_{moy}}{E_{P}} = \frac{1000W}{0.35 J}  = 2857 Hz \approx 3KHz}
    \end{equation}
    
\end{itemize}

\subsection{Calculation of $\Delta V$ and Deorbiting Dynamics}
The calculation of the required velocity change $\Delta V$ and the transferred impulse is based on three main elements: \\
\begin{itemize}
    \item \textbf{The momentum coupling coefficient ($C_{m}$)}\\
    When a LASER beam interacts with the surface of a material at a fluence above its ablation threshold (aluminium in this case),
    the material vaporises instantaneously. This ejection of matter, the \textbf{rocket effect}, generates a recoil force. 
    This coefficient measures the efficiency with which the LASER transfers momentum to the debris. The coupling is defined by the following fundamental relationship \cite{Phipps2012}:\\
    \begin{equation}
        \boxed{C_{m} = \frac{I}{E_{LASER}} = \frac{\Delta p}{E_{LASER}}}
    \end{equation}
    with \textbf{I} the total generated impulse (\textbf{$N.S$}) and \textbf{$E_{LASER}$} the total absorbed LASER energy in (\textbf{J}). $C_{m}$ is expressed in (\textbf{$N.s/J$}).\\
    According to the experimental data in \cite{Phipps2012} for a nanosecond LASER at a wavelength of $\lambda = 1064nm$, the stabilised value at peak fluence is equal to $C_{m} = 30 \mu N.s/J$\\

    \item \textbf{Average generated force $F_{moy}$}\\
    This force is therefore given by the expression below:
    \begin{equation}
       \boxed{ F_{moy} = c_{m} * P_{moy} = 30.10^{-6} * 1000 = 0.03N = 3 mN}
    \end{equation}\\

    \item \textbf{Orbital dynamics }\\
    In our case, we opt for a staged atmospheric re-entry, lowering the debris perigee and allowing the atmosphere to complete the process. We therefore apply our LASER impulse directly opposite to the velocity vector \textbf{$\vec{v_{debris}}$} of the debris at apogee.\\
    According to the vis-viva equation and Kepler's laws, we define the circular velocity \textbf{$\vec{{v_{c}}}$} of the debris in its initial orbit (at \textbf{$350 Km$} ).
    This is the velocity required for the centrifugal force to balance Earth's gravity exactly and maintain its trajectory.
    \begin{equation}
       \boxed{ {v_{c}} = \sqrt{\frac{\mu}{R_{E} + h}} = \sqrt{\frac{3.986.10^{14}}{6371000 + 350000}} = 7701 m/s }
    \end{equation}

To lower its perigee to a critical altitude of \textbf{$h_{p}= 100Km$}, the velocity at the new apogee \textbf{$\vec {v_{a}}$} is given by: 
\begin{equation}
  \boxed{  v_{a} = \sqrt {\mu \left(\frac{2}{r_{a}} -\frac{2}{r_{a} + r_{P}}\right )} = \sqrt {3.986.10^{14}\left(\frac{2}{6371000 + 350000} -\frac{2}{6721000 + 6471000}\right )} = 7628 m/s}
\end{equation}
Here, $r_{a} = R_{E} + h$, represents the apogee radius of the new trajectory;\\
$r_{P} = R_{E} + h_{P}$, represents the perigee radius of the new trajectory.\\
We select a LEO orbit at $h = 350Km$, which is the initial orbit of our debris (corresponding to the initial orbit of our CubeSat).\\
$h_{P} = 100Km$ corresponds to the target altitude of the debris. At this orbit, the atmosphere is sufficiently dense to decelerate our debris and cause re-entry. \\
$\mu = 3.986.10^{14} m^{3}/s^{2}$ is Earth's standard gravitational parameter.\\
$R_{E} \approx 6371 Km$ represents Earth's mean radius.\\
The expression for $\Delta V$ is therefore written as follows: 
\begin{equation}
    \boxed{\Delta V = v_{c} - v_{a} = 7701 -7628 = 73m/s}
\end{equation}
Numerical result: $\boxed{\Delta V = 73  m/s}$.\\
The total required firing duration can also be calculated. \\
\begin{equation}
    \boxed{\Delta t = \frac{m_{debris} * \Delta V}{F_{moy}} = \frac{1* 73}{0.03} = 2433 s }
\end{equation}
The required firing duration $\Delta t $ is approximately $40$ minutes and $33$ seconds to impart a velocity change of $\Delta V = 73$ $m/s$ to our debris,
divert it from its initial orbit and cause atmospheric re-entry. \cite{Vallado2013}\\

But it's impossible to keep the contact between the debris and the CubeSat for 40 minutes.
It's no nessecary to create a big $\Delta V$ if you change the orbit with a small $\Delta V$ and wait for the debris to re-enter the atmosphere it's a success.
Thanks to this experimental project, we will be able to validate the feasibility of our concept and the effectiveness of our LASER ablation system in space.

\end{itemize}
